\section{Introduction}\label{sec:introduction}

Let \(V=(P,Q)\) be a real polynomial vector field on \(\RR^2\).
A periodic orbit is the image of a nonconstant periodic solution of
\[
 \dot x=P(x,y),\qquad \dot y=Q(x,y).
\]
It is a \emph{limit cycle} if some open neighborhood of that image contains
no other periodic orbit. We count geometric images once, without
multiplicity.

\begin{theorem}[Uniform boundedness]\label{thm:uniform}
For each integer \(d\geq1\), there is a finite nonnegative integer \(B(d)\)
such that every real planar polynomial vector field of degree at most \(d\)
has at most \(B(d)\) limit cycles in the whole plane.
\end{theorem}

There is no restriction on the coefficients, the locations of the cycles,
or their stability. Theorem~\ref{thm:uniform} gives a positive resolution of
the uniform boundedness assertion in the second part of Hilbert's
sixteenth problem. It does not furnish an effective formula for \(B(d)\).

\subsection{Historical context}

Hilbert's sixteenth problem asks, in its differential-equation part, for
the maximum number and arrangement of limit cycles of planar polynomial
systems \citep[Problem~16, p.~465]{Hilbert1902}. Two finiteness assertions
must be distinguished. Individual finiteness concerns one fixed vector
field; uniform boundedness asks for a bound depending only on its degree
\citep[pp.~301--304]{Ilyashenko2002}. Even finiteness for every member of a
family does not by itself bound the numbers occurring as its coefficients
vary.

Dulac claimed a proof of individual finiteness in 1923
\citep{Dulac1923}; after a gap in that argument was identified
\citep{Ilyashenko1985}, \'Ecalle and Ilyashenko developed proofs using
complex and asymptotic analysis of return maps
\citep{Ilyashenko1990,Ilyashenko1991,Ecalle1992}. Yeung identifies a
coefficient-closure obstruction in a leading-term argument of
Ilyashenko's 1991 monograph \citep{Yeung2025}. This concerns that proof
mechanism; it is not a counterexample to individual finiteness.

The geometric approach to uniformity studies the limiting sets of
periodic orbits and bounds the number of nearby cycles in a parameter
family. A foundational local result is Bautin's sharp bound of three
for cycles bifurcating from a nondegenerate weak focus or center of a
quadratic vector field, under perturbations within the quadratic family
\citep[\S3]{Bautin1952}.
Roussarie develops the study of limiting periodic sets through graphics
and desingularization \citep{Roussarie1998}.
Ilyashenko and Yakovenko prove finite cyclicity for elementary
polycycles in generic finite-parameter families
\citep{IlyashenkoYakovenko1995}; Kaloshin subsequently gives an explicit
cyclicity bound in terms of the number of parameters \citep{Kaloshin2003}.
The elementary hypothesis requires each singularity of the polycycle
to have at least one nonzero eigenvalue. It and the genericity assumptions
restrict these results. Another major partial advance concerns the
infinitesimal problem: Binyamini, Novikov, and Yakovenko give explicit
bounds for zeros of Abelian integrals and for cycles arising from
nonsingular Hamiltonian ovals in the stated nonconservative perturbations
\citep{BinyaminiNovikovYakovenko2010}. These results illustrate both the
power of local transition analysis and the extra work required at
unrestricted degenerations.

There is a parallel development through definability and quasianalytic
asymptotics. Kaiser, Rolin, and Speissegger construct an o-minimal
expansion containing the transition maps at nonresonant hyperbolic
singularities. For small analytic unfoldings that preserve the singular
points and their linear parts, they obtain uniform bounds on isolated
return-map fixed points near polycycles with such singularities
\citep{KaiserRolinSpeissegger2009}.
Galal, Kaiser, and Speissegger construct Ilyashenko algebras with
injective transserial asymptotics and a corresponding Hardy field
\citep{GalalKaiserSpeissegger2020}. They prove closure under field
operations and differentiation for germs in one variable. The
multivariable closure needed for a parameter argument is a separate
issue. The present
proof establishes separation, differentiated implicit closure, and
projection bounds for its own passage class. Those properties are proved
in Sections~\ref{sec:separation}--\ref{sec:elimination} and
\ref{sec:counting}; definability of arbitrary passage maps is not an
assumption.

The proof below treats coefficient variation and degeneration together.
It does not use individual finiteness as a premise. This is also why
cycles that move arbitrarily far from the origin and nonhyperbolic
isolated cycles must be included in the last step of the argument.

\subsection{Proof strategy}

Choose a short transverse interval through a point of a periodic orbit,
with coordinate \(r\). Following nearby trajectories once around the
orbit defines a real analytic return map \(\Pi(r)\). The nearby limit
cycles correspond to isolated zeros of the displacement \(\Pi(r)-r\).
A simple zero has \(\Pi'(r)\ne1\); such a cycle is hyperbolic. The
difficulty is to bound these zeros while the coefficients vary and the
passages making up the return map degenerate.

A local saddle already exhibits the scales that occur. For
\(\dot x=x\), \(\dot y=-\lambda y\), with \(\lambda>0\), the
passage from \((r,1)\) to \((1,r^\lambda)\), \(0<r<1\), has
flight time \(L=\log(1/r)\) and contraction exponent \(W=\lambda L\).
When \(r\to0\) and \(\lambda\to0\), the first clock can diverge
while the second stays bounded or diverges at a different rate. The
later passage models treat such differing scales together with their
smaller correction terms.

We represent a return by equations matching the endpoints of finitely
many local passages. This formulation has a predecessor in Ilyashenko
and Yakovenko's cyclic systems for elementary polycycles
\citep[\S0.4]{IlyashenkoYakovenko1995}. Here the cuts are allowed to
move along the orbit, giving continuous families of representations.
Section~\ref{sec:counting} shows that, within each regular chart of a
fixed matching system, representations of different hyperbolic cycles
lie in different connected components of the fiber. It therefore
remains to bound those components uniformly as
the field coefficients and preparation parameters vary.

The counting argument reduces this uniform bound to a different
assertion: each fixed finite auxiliary system has only finitely many
isolated solutions when all its variables, including those later used
as parameters, are allowed to vary. We call this \emph{absolute
isolated-zero finiteness}. It must hold for the differential and
algebraic systems created by the counting argument, including their
graph relations and multiplier equations, with the original open
domains retained. The analytic and geometric parts of the proof
establish these hypotheses for the matching systems.

Here is the organization of the proof.

\begin{enumerate}[label=\arabic*.]
\item
Sections~\ref{sec:flags}--\ref{sec:elimination} prove the abstract
absolute-zero criterion in Theorem~\ref{thm:absolute-zero}. Along an alleged escaping sequence, a finite
saturated flag records the ordered large scales and their exact
logarithmic relations. Theorem~\ref{thm:separation} compares two lateral trees
of asymptotic coefficients and identifies the first nonzero term. If
both trees vanish, the function vanishes exactly. A differentiated
implicit calculus and an induction on ordinary variables then turn this
information into a local chart through each sufficiently far solution.
An independent large coordinate survives in a zero chart, contradicting
isolation. The differentiated implicit calculus is stated in
Theorem~\ref{thm:calculus}.

\item
Sections~\ref{sec:additive}--\ref{sec:mixed} verify the packet and
retestability hypotheses of that criterion for the analytic passage maps
used later
(Theorems~\ref{thm:additive-packets}, \ref{thm:regular-packets},
and~\ref{thm:mixed-packets}). The models
include additive transfers, regular or stable transfers with one large
clock, and transfers with two clocks that may be comparable or widely
separated. Their arguments are treated independently
before endpoint relations are imposed. In particular, exponentially
small terms and the differences between lateral determinations remain
part of the packet data.

\item
Sections~\ref{sec:subdivision} and~\ref{sec:terminal} give a finite
geometric description for fields of bounded degree. Strong
one-variable subanalytic preparation is used on the scalar fields,
before passage maps are introduced. It supplies finitely many boxes
and annuli (Theorem~\ref{thm:finite-templates}), which reduce to the
preceding analytic models (Theorem~\ref{thm:terminal-reduction}).
A crossing argument for periodic Jordan curves bounds the length of
the resulting passage words. This finite description depends on the
degree; it is fixed before any sequence is chosen for an absolute-zero
test.

\item
Section~\ref{sec:counting} constructs the matching equations with their
exact graph relations and verifies absolute finiteness for the required
closure (Theorem~\ref{thm:matching-system}). The independent counting
principle is Theorem~\ref{thm:projection-count}.
Within each regular chart, representations of different hyperbolic
cycles lie in different connected components of the fiber.
A suitable small signed rotation produces at least as many hyperbolic
cycles as any prescribed finite collection of isolated cycles
(Lemma~\ref{lem:counting-rotation}). Spatial rescaling then gives the
assertion throughout the plane.
\end{enumerate}

Figure~\ref{fig:proof-map} records where the analytic and geometric
parts meet in the counting argument.
\begin{figure}[tbp]
\centering
\begin{tikzpicture}[
  box/.style={draw,rounded corners=2pt,align=center,
    text width=3.8cm,minimum height=1.28cm,inner sep=5pt,
    font=\small},
  wide/.style={box,text width=12.7cm},
  half/.style={box,text width=5.8cm},
  flow/.style={-{Stealth[length=2mm]},semithick},
  every node/.style={outer sep=2pt}]
\node[box] (criterion) at (-4.6,0)
  {Abstract isolated-zero criterion\\
   Sections~\ref{sec:flags}--\ref{sec:elimination}};
\node[box] (packets) at (0,0)
  {Verified passage packets\\
   Sections~\ref{sec:additive}--\ref{sec:mixed}};
\node[box] (geometry) at (4.6,0)
  {Finite geometric templates; reductions to transfer models\\
   Sections~\ref{sec:subdivision}--\ref{sec:terminal}};
\node[wide] (matching) at (0,-2.3)
  {Matching systems: absolute isolated-zero finiteness\\
   for every finite system in the required differential and algebraic closure\\
   Theorem~\ref{thm:matching-system}};
\node[half] (projection) at (-3.5,-4.45)
  {Projection count for regular fibers\\
   Theorem~\ref{thm:projection-count}};
\node[half] (hyperbolic) at (3.5,-4.45)
  {Uniform hyperbolic-cycle bound\\in a fixed square};
\node[wide] (cycles) at (0,-6.5)
  {A bound depending only on degree, throughout the plane\\
   Theorem~\ref{thm:uniform}};
\draw[flow] (packets.north) -- ++(0,.55)
  -| node[pos=.25,above,font=\small] {terminal reductions} (geometry.north);
\draw[flow] (criterion.south)--(matching.north -| criterion.south);
\draw[flow] (packets)--(matching);
\draw[flow] (geometry.south)--(matching.north -| geometry.south);
\draw[flow] (matching.south -| hyperbolic.north)--(hyperbolic.north);
\draw[flow] (projection)--(hyperbolic);
\draw[flow] (hyperbolic.south)--
  node[right,align=left,font=\small] {rotation and\\rescaling}
  (cycles.north -| hyperbolic.south);
\end{tikzpicture}
\caption{Main theorem interfaces. Matching combines the finite geometric
 description with the packet criterion, retaining the exact graph ties
 and open domains. The independent projection theorem converts absolute
 finiteness for the full required closure into uniform component bounds.
 The sequence-dependent flags and internal factorizations used in the
 absolute-zero test do not enlarge the finite geometric alphabet.
 The packet input to Sections~\ref{sec:subdivision}--\ref{sec:terminal}
 concerns the terminal reductions, not the preparation of scalar fields.}
\label{fig:proof-map}
\end{figure}


The analytic flags and charts in an absolute-zero contradiction may
depend on the chosen sequence. They are not used as a finite covering
of coefficient space. The finite geometric description and the
projection-counting argument supply the uniformity.

For a first reading, the theorem statements cited above and the final
proof in Section~\ref{sec:counting} display the chain of implications.
Definition~\ref{def:counting-domain} specifies the precise closure
required at the counting interface. The preceding sections establish
its analytic and geometric inputs.

\subsection{Analytic ingredients}

Three features of the argument deserve emphasis.

A \emph{packet} records an actual analytic function together with two
trees of successive asymptotic expansions and the domain, transition,
and error estimates relating them. The two signs refer to lateral
determinations on either side of the real path; higher-band raw
realizations may be smooth interpolations of holomorphic determinations.

First, the separation theorem uses a \emph{tree} of expansions on
successive complex bands. It does not assume a simultaneous convergent
series in all small exponentials. The hypotheses record retained
arguments, permissible error costs, differentiated transition
estimates, and growth on the full fringes of the domains.
The proof upgrades arbitrarily high fixed exponential decay to exact
vanishing by weighted Cauchy corrections and maximum principles.
The domain and error conditions are indispensable: real asymptotic
flatness alone gives no such conclusion.

Second, regular ordinary implicit changes must be compatible with
dependent scales that can rotate rapidly under tiny complex
perturbations. The construction uses high real Taylor jets and
interpolation into small comparison tubes, with a fixed finite
coherence order at each stage. It does not require one ordinary
complex neighborhood valid for all derivative orders.
The later center-lifting argument additionally keeps the power
exponents controlling maps, inverses, and neighborhoods fixed while
the accuracy of auxiliary records is increased. This quantitative
property permits exact correction of approximate chart hits.

Third, the projection argument is formulated for an analytic class
closed under the particular derivatives, graph variables, and
multiplier equations used in its proof. Proper barriers and a fixed
generic tilt reduce the number of parameters. The resulting
absolute-to-uniform implication is independent of the particular
planar passage models and can be applied to other classes for which
the same hypotheses are established.

\subsection{Notation and conventions}

The technical constructions are local at a limiting ordinary analytic
germ. An asymptotic assertion is made for a fixed finite request:
a finite collection of expressions, derivative orders, Taylor
accuracies, and transition estimates. Constants, sufficiently far
starting points, and smaller neighborhoods may depend on that
request. Whenever an exponent must be independent of a later
accuracy choice, this is stated explicitly.

The notation used throughout the analytic part is summarized below;
see Definitions~\ref{def:flags-flag} and~\ref{def:packet}.
Local clock and state variables are introduced separately in each
passage model.

\begin{center}
\begin{tabular}{@{}ll@{}}
\toprule
Notation & Meaning\\
\midrule
\(Z_1\gg\cdots\gg Z_m\gg1\) & Ordered scales of a finite flag\\
\(f_i=\log Z_i,\quad h_i=e^{-Z_i}\) & Logarithmic scales and small exponentials\\
\(S\) & Tuple of independent, or free, flag scales\\
\(t,t_c\) & Ordinary inputs and their interior limiting germ point\\
\(b_i\) & Bounded ordinary backgrounds in exact log relations\\
\(\sigma(i)\) & Leading later index in a dependent log relation\\
\(s=u+iv\) & Complex parameter of a selected flag path\\
\(\Lambda_i=f_i'\) & Logarithmic rate along the real path\\
\(V_i=e^{f_i}/\Lambda_i\) & Precision scale at level \(i\)\\
\(S^{j,\pm}_\alpha\) & Raw packet in band \(j\) at prefix \(\alpha\)\\
\(h^\nu=\exp(-\sum_i\nu_iZ_i)\) & Exponential shift with lexicographic order\\
\bottomrule
\end{tabular}
\end{center}

Lexicographic order starts at the largest scale. A positive exponent
vector has positive first nonzero coordinate; later coordinates need
not be nonnegative. A vanishing coefficient means a vanishing
ordinary analytic germ. Complex logarithms and powers always use the
specified lifted determinations. A dependent scale is evaluated at
its \emph{current} background values, including after an implicit
change of ordinary variables.

All differentiations of primitive packets are in their independent
arguments. Equations tying those arguments are imposed only
afterward; differentiation of the resulting physical maps uses the
chain rule on the tied graph. This distinction is used repeatedly in
the terminal reductions and the matching equations.
